\documentclass{article}
\usepackage{amsmath, amssymb}
\usepackage{amsthm}

\author{Michael Naumov}
\date{}
\setlength{\parindent}{0pt}
\setlength{\parskip}{0.5em}


\title{1.2.1\mbox{-}13}

\begin{document}
\maketitle

\( A1: m > 0, n > 0, T = 0 \)

\( A2: c = m > 0, d = n > 0, a = b' = 0, a' = b = 1, T = 1 \)

\( A3: am + bn = d, a'm + b'n = c = qd + r, 0 \leq r < d, \gcd(c, d) = \gcd(m, n), T \leq 3(n-d) + 2 \)

\( A4: am + bn = d = \gcd(m, n), d > 0, T \leq 3(n - d) + 3 \)

\( A5: am + bn = d, a'm + b'n = c = qd + r, 0 < r < d, \gcd(c, d) = \gcd(m, n), T \leq 3(n-d) + 3 \)

\( A6: am + bn = d, a'm + b'n = c, d > 0, \gcd(c, d) = \gcd(m, n), T \leq 3(n-d) + 1 \)

Let's prove the additional statements for \( T \) by induction.

\( A1: T = 0 \), by definition

\( A2: T=1 \), by definition

\( A3: T \leq 3(n-d)+2 \)

After \( E2: d=n, T_{A3}=T_{A2} + 1 = 2 \leq 3(n-d) + 2 \)

After \( E4: T_{A3} = T{A6} + 1 \leq 3(n-d)+1+1 = 3(n-d)2 + 2 \)

\( \implies A3 \)

\( A4: d > 0, T \leq 3(n-d)+3 \)

Part \( d > 0 \) always holds, as \( \gcd \) cannot be \( 0 \).

\( T_{A4} = T_{A3} + 1 \leq 3(n - d)+2 + 1 = 3(n-d)+3 \implies A4 \)

\( A5: T \leq 3(n-d) + 3 \)

\( T_{A5} = T_{A3} + 1 \leq 3(n - d)+2 + 1 = 3(n-d)+3 \implies A5 \)

\( A6: T \leq 3(n-d)+1 \)

\( E4 \) assigns \( d \leftarrow r \). Let's \( d_{0} \) be value of \( d \) before assigning.

\( T_{A6} = T_{A5} + 1 \leq 3(n-d_{0})+3+1 = 3(n-d_{0})+4 = 3(n-r)+1+3(r-d_{0}+1) \)

\( 0 \leq r \leq d_{0} - 1 \implies r + 1 - d_{0} \leq 0 \implies T_{A6} \leq 3(n-r)+1 = 3(n-d)+1 \implies A6 \)


\end{document}
